\documentclass{beamer}
\usepackage{amssymb}

\begin{document}
  
  \begin{frame}
    \textit{Definición}: Dados $A, B$ conjuntos, decimos que 
    $A$ es quipotente a $B$ si existe $f: A \to B$ biyectiva y lo
    simbolizamos por $A \equiv B$.
    \newline\newline
    \pause 

    \textit{Teorema}: Sean $A, B$ conjuntos arbitrarios, entonces
    \begin{itemize}[<+->]
      \item $A \equiv A$.
      \item $A \equiv B \iff B \equiv A$.
      \item $(A \equiv B \land B \equiv C) \rightarrow A \equiv C$.
      \item $P(A) \equiv\, ^2 A, 2 \in \mathbb{N}, 2 = \{0, 1\}$.
    \end{itemize}
  \end{frame}

  \begin{frame}
    \textit{Teorema}: $A \equiv B \rightarrow P(A) \equiv P(B)$.
    \newline\newline
    \pause

    \textit{Teorema}: Si $A \equiv D$ y $B \equiv E$, entonces 
    \begin{itemize}[<+->]
      \item Si $A \cap B = \emptyset = D \cap E$, entonces $A \cup B \equiv D \cup E$.
      \item $A \times B \equiv D \times E$.
      \item $^B A \equiv\, ^E D$.
    \end{itemize}
  \end{frame}

  \begin{frame}
    \textit{Teorema}: Sean $A, B, C$ conjuntos arbitrarios 

    \begin{itemize}[<+->]
      \item Si $A \cap B = \emptyset$ entonces $^{A \cup B} C \equiv\, ^A C \times\, ^B C$.
      \item $^A (B \times C) \equiv\, ^A B \times\, ^A C$.
      \item $^A (^B C) \equiv\, ^{A \times B} C$.
    \end{itemize}
  \end{frame}

  \begin{frame}
    \textit{Definición}: Decimos que $A$ es dominado por $B$ si existe 
    una función inyectiva de $A$ hacia $B$ y se denota por 
    $A \curlyeqprec B$, si no existe una unción biyectiva de $A$ en 
    $B$, decimos que la dominancia es estricta y se denota por $A \prec B$.
    \newline\newline
    \pause
    
    \textit{Teorema}: Sean $A, B, C$ conjuntos: 
    \pause
    \begin{itemize}[<+->]
      \item $A \curlyeqprec A$.
      \item $A \curlyeqprec P(A)$.
      \item Si $A \subseteq B$, entonces $A \curlyeqprec B$.
      \item Si $A \curlyeqprec B$ y $B \curlyeqprec C$ entonces $A \curlyeqprec C$.
      \item Si $A \curlyeqprec B$ y $B \equiv C$ entonces $A \curlyeqprec C$.
      \item Si $A \equiv B$ y $B \curlyeqprec C$ entonces $A \curlyeqprec C$.
      \item $A \curlyeqprec B$ si y sólo si existe $D \subseteq B$ tal que $A \equiv D$.
      \item $A \equiv B \rightarrow (A \curlyeqprec B \land B \curlyeqprec A)$.
    \end{itemize}
  \end{frame}

  \begin{frame}
    \textit{Teorema (Cantor)}: Sea $A$ cualquier conjunto. No existe una 
    función sobreyectiva de $A$ hacia $P(A)$.
    \newline\newline
    \pause
    \textit{Teorema}: Sea $A$ un conjunto y $F: P(A) \to P(A)$ una función
    creciente:
    \pause

    \begin{itemize}[<+->]
      \item $T = \bigcap \{x \subseteq A: F(x) \subseteq x\}$ entonces $F(T) = T$.
      \item $T' = \bigcup \{x \subseteq A: x \subseteq F(x)\}$ entonces $F(T') = T'$.
      \item Si $F(y) = y$ entonces $T \subseteq y  \subseteq T'$.
    \end{itemize}
  \end{frame}
\end{document}
